\documentclass[10pt,a4paper]{amsart}
\usepackage[margin=19mm]{geometry}
\usepackage{amsmath,amssymb,mathtools}
\usepackage[scr=rsfs,cal=euler]{mathalfa}
\usepackage{xfrac}
\usepackage[unicode,hidelinks,bookmarks=false]{hyperref}
\newcommand{\ie}{i.e.\ }
\newcommand{\cf}{cf.\ }
\newcommand{\resp}{respectively\ }
\newcommand{\Subsection}{\subsection}
\providecommand{\nobreakdash}{\nobreak\hbox{-}}
\setlength{\emergencystretch}{2em}
\multlinegap0pt
\usepackage{amssymb}
\newtheorem{theorem}{Theorem}
\newtheorem{lemma}{Lemma}
\newtheorem{proposition}{Proposition}
\newcommand{\Z}{\mathbb{Z}}
\title{Proofs for Andrews' Conjectures 5 and 6 on $v_1(q)$}
\author{Mohamed El Bachraoui}
\date{Source: arXiv:2604.08013v1 (CC BY 4.0)}
\begin{document}
\maketitle

\noindent\emph{Source and licence.} Proofs for Andrews' Conjectures 5 and 6 on $v_1(q)$. Version arXiv:2604.08013v1 (CC BY 4.0), distributed by arXiv under the Creative Commons Attribution 4.0 International licence, http://creativecommons.org/licenses/by/4.0/, as recorded in the arXiv licence metadata for this exact version and shown on the abstract page https://arxiv.org/abs/2604.08013v1. Full licence notice: \texttt{licenses/arxiv-2604.08013-CC-BY-4.0.txt}.\par\smallskip

\noindent\textit{Editorial scope.} Counted unit: the article's theorem thm:main (source0.tex lines 94--108). The article's complete original proof of it is delivered with it (source0.tex lines 948--1015, one \texttt{proof} environment of 68 lines, which the article places away from the statement). No mathematics is written, repaired or re-derived: the statement and the proof are the article's own lines, in the article's own notation, and no retained line is substituted or re-flowed.  Retained uncounted as the source-local context the proof needs: lines 64--66; lines 241--255; lines 257--271; lines 281--283; lines 913--940.  Nothing was omitted from any retained span, so the package carries no omission record.  No retained line contains a citation, so the wrapper carries no bibliography and no bibliography entry.  No retained line contains a figure environment, an image-inclusion command or drawing code, so no image is delivered and the package declares no asset dependency.  Editorial formatting only, in addition: an AMS \texttt{amsart} preamble that restates the article's own declarations and macros used by the retained lines, namely the environments \texttt{theorem}, \texttt{lemma}, \texttt{proposition} on separate counters, together with the standard packages those lines need.  The retained environments print in this document as Theorem 1, Lemma 1, Proposition 1 in order of appearance, whereas the article prints the same items with the numbering of its own full document.  The complete native source of the article as distributed in the arXiv e-print for this exact version is bundled unchanged as \texttt{sources/198131/source0.tex}.
\begin{equation}\label{eq:v1-def}
  v_1(q):=\sum_{n\ge 0}\frac{q^{n(n+1)/2}}{(-q^2;q^2)_n}=\sum_{n\ge 0}V_1(n)q^n.
\end{equation}

\begin{theorem}\label{thm:minus-family}
Let $m\in \mathcal M_-$, and set
\[
  N_m:=2\left\lfloor \frac{\alpha(m+1/4)^2}{2}\right\rfloor.
\]
Then $N_m$ is the unique even integer satisfying
\begin{equation}\label{eq:Nm-def}
  N_m<\alpha\left(m+\frac14\right)^2<N_m+2.
\end{equation}
For all sufficiently large $m\in \mathcal M_-$,
\[
  V_1(N_m),\qquad V_1(N_m+1),\qquad V_1(N_m+2)
\]
have the same sign. Moreover, both $N_m$ and $N_m+2$ are strict local minima of the sequence $|V_1(n)|$.
\end{theorem}

\begin{theorem}\label{thm:plus-family}
Let $m\in \mathcal M_+$, and set
\[
  M_m:=2\left\lfloor \frac{\alpha(m+3/4)^2-1}{2}\right\rfloor+1.
\]
Then $M_m$ is the unique odd integer satisfying
\[
  M_m<\alpha\left(m+\frac34\right)^2<M_m+2.
\]
For all sufficiently large $m\in \mathcal M_+$, the three numbers
\[
  V_1(M_m),\qquad V_1(M_m+1),\qquad V_1(M_m+2)
\]
have the same sign. Moreover, both $M_m$ and $M_m+2$ are strict local minima of the sequence $|V_1(n)|$.
\end{theorem}

\begin{lemma}\label{lem:alpha-not-Z}
We have $\alpha\notin \Z$.
\end{lemma}

\begin{proposition}\label{prop:initial-values}
A direct expansion of \eqref{eq:v1-def} gives
\[
\begin{aligned}
V_1(292)&=-367,\qquad V_1(293)=4,\qquad V_1(294)=375,\\
V_1(295)&=9,\qquad V_1(296)=-381,
\end{aligned}
\]
\[
\begin{aligned}
V_1(409)&=-465,\qquad V_1(410)=27,\qquad V_1(411)=473,\\
V_1(412)&=4,\qquad V_1(413)=-497,
\end{aligned}
\]
\[
\begin{aligned}
V_1(544)&=6195,\qquad V_1(545)=-18,\qquad V_1(546)=-6309,\\
V_1(547)&=-20,\qquad V_1(548)=6418,
\end{aligned}
\]
\[
\begin{aligned}
V_1(701)&=8365,\qquad V_1(702)=-273,\qquad V_1(703)=-8550,\\
V_1(704)&=-224,\qquad V_1(705)=8716.
\end{aligned}
\]
In particular, the initial values $293$, $410$, $545$, and $702$ satisfy Andrews' Conjectures~5 and~6.
\end{proposition}

\begin{theorem}\label{thm:main}
There exists a sequence of integers $(N_n)_{n\ge 5}$ satisfying
\[
  N_5=293,\qquad N_6=410,\qquad N_7=545,\qquad N_8=702,\qquad N_n>10n^2\quad (n\ge 9),
\]
such that for every $n\ge 5$ the three numbers
\[
  V_1(N_n),\quad V_1(N_n+1),\quad V_1(N_n+2)
\]
have the same sign, and one of
\[
  |V_1(N_n)|,\quad |V_1(N_n+1)|,\quad |V_1(N_n+2)|
\]
is a local minimum of the sequence $|V_1(j)|$.
\end{theorem}

\begin{proof}[Proof of Theorem~\ref{thm:main}]
By Proposition~\ref{prop:initial-values}, the integers
\[
  293,\qquad 410,\qquad 545,\qquad 702
\]
already satisfy the required same-sign and local-minimum properties.

Next, by definition,
\[
  N_m=\alpha\left(m+\frac14\right)^2+O(1),
  \qquad
  M_m=\alpha\left(m+\frac34\right)^2+O(1).
\]
The proof of Lemma~\ref{lem:alpha-not-Z} showed that $\alpha>10$. Hence
\[
  \frac{N_m}{(m+8)^2}\to \alpha,
  \qquad
  \frac{M_m}{(m+8)^2}\to \alpha,
\]
so after discarding finitely many elements from $\mathcal M_-$ and $\mathcal M_+$ we may assume simultaneously that Theorems~\ref{thm:minus-family} and \ref{thm:plus-family} apply and that
\[
  N_m>10(m+8)^2\quad(m\in\mathcal M_-),
  \qquad
  M_m>10(m+8)^2\quad(m\in\mathcal M_+).
\]
Now choose recursively $m_j\in\mathcal M_-$ and $\ell_j\in\mathcal M_+$ so that
\[
  m_j\ge 2j-1,
  \qquad
  \ell_j\ge 2j,
\]
and, with
\[
  L_{2j-1}:=N_{m_j},
  \qquad
  L_{2j}:=M_{\ell_j},
\]
we have
\[
  702<L_1<L_2<L_3<\cdots.
\]
This is possible because after each finite exclusion both admissible families still contain arbitrarily large elements. For these choices,
\[
  L_{2j-1}>10(m_j+8)^2\ge 10(2j+7)^2,
  \qquad
  L_{2j}>10(\ell_j+8)^2\ge 10(2j+8)^2.
\]
Hence $L_j>10(j+8)^2$ for every $j\ge1$.

Finally define
\[
  N_5:=293,\qquad N_6:=410,\qquad N_7:=545,\qquad N_8:=702,
\]
and for $j\ge1$ put
\[
  N_{8+j}:=L_j.
\]
Then $N_n>10n^2$ for every $n\ge 9$. Moreover, for each $n\ge9$, the three coefficients
\[
  V_1(N_n),\qquad V_1(N_n+1),\qquad V_1(N_n+2)
\]
have the same sign, and one of the numbers
\[
  |V_1(N_n)|,\qquad |V_1(N_n+1)|,\qquad |V_1(N_n+2)|
\]
is a local minimum of the sequence $|V_1(j)|$, by Theorem~\ref{thm:minus-family} or Theorem~\ref{thm:plus-family} according as $N_n=L_j$ with $j$ odd or even. 
Together with the four initial values, this proves the theorem, i.e. Andrews' Conjectures~5 and~6.
\end{proof}

\end{document}
